% Source: https://artofproblemsolving.com/community/c3h1232121_3d_geometry
% Stable identity: aops-topic-1232121-q01
The centres form a regular tetrahedron $T$ of edge $2$. The faces of the big tetrahedron are parallel to the faces of $T$, each at distance $1$ (outside $T$), because each face is tangent to three spheres. So both tetrahedra have the same centre, and the inradius grows by $1$.

A regular tetrahedron of edge $e$ has height $e\sqrt{\dfrac{2}{3}}$ and inradius one quarter of the height, $\dfrac{e}{2\sqrt6}$. For $T$ this is $\dfrac{1}{\sqrt6}$, so the big tetrahedron has inradius $1+\dfrac{1}{\sqrt6}$ and edge
\[2\sqrt6\left(1+\dfrac{1}{\sqrt6}\right)=\boxed{2+2\sqrt6}.\]
In the side view below (the plane through one edge and the midpoint of the opposite edge), the tetrahedron of centres (dashed) sits inside the outer one, every face at distance $1$.

\begin{asy}
usepackage("amsmath");
size(4cm);
defaultpen(fontsize(10pt));
real e=2; real E2=2+2*sqrt(6);
pair f(real ed, pair c){ return c; }
real h1=e*sqrt(2/3), h2=E2*sqrt(2/3);
real m1=e*sqrt(3)/2, m2=E2*sqrt(3)/2;
// side view: triangle with base = median of a face, apex = top vertex
pair O=(0,0);
pair B1=(-m1/3,-h1/4), A1=(2*m1/3,-h1/4), T1=(0,3*h1/4);
pair B2=(-m2/3,-h2/4), A2=(2*m2/3,-h2/4), T2=(0,3*h2/4);
draw(B2--A2--T2--cycle,linewidth(1)); draw(B1--A1--T1--cycle,blue+dashed);
dot(A1);dot(T1); draw(circle(T1,1),gray); draw(circle(A1,1),gray);
label("outer tetrahedron",T2,N);
label("\mbox{}",(min(currentpicture,user=true)+max(currentpicture,user=true))/2,black); // keeps the text after the figure black when LaTeX typesets the labels inline
\end{asy}
